\section{Number of Enclaves}
\label{sec:graph-number-of-enclaves}

\problemstatement{Given a binary grid where $0$ is sea and $1$ is
land, a move walks $4$-directionally from one land cell to another.
Count the number of land cells from which it is \textbf{impossible}
to walk off the boundary of the grid in any number of moves.}

\begin{patternbox}
This is Section~\ref{sec:graph-surrounded-regions}'s boundary-first
elimination technique again, with the roles renamed: land cells
connected to the border can ``escape,'' so traverse from every
\textbf{border land cell} first, marking everything reachable as
escapable; every \textbf{unmarked} land cell remaining afterward is an
enclave, and the answer is simply the count of those.
\end{patternbox}

\subsection*{Algorithm}

\begin{enumerate}
  \item For every land cell on the border, DFS (or BFS) outward,
        marking every $4$-directionally connected land cell as visited
        (mutate the grid in place, setting visited land to $0$, exactly
        as the previous section reused the grid itself as the
        \textit{visited} structure).
  \item Scan the entire grid and count every remaining land cell
        ($1$) -- these are cells no border-seeded traversal ever
        reached, i.e.\ the enclaves.
\end{enumerate}

\subsection*{C++ implementation}

\begin{lstlisting}
#include <bits/stdc++.h>
using namespace std;

void markEscapable(int r, int c, vector<vector<int>>& grid) {
    int rows = grid.size(), cols = grid[0].size();
    if (r < 0 || r >= rows || c < 0 || c >= cols) return;
    if (grid[r][c] != 1) return;

    grid[r][c] = 0;
    markEscapable(r - 1, c, grid);
    markEscapable(r + 1, c, grid);
    markEscapable(r, c - 1, grid);
    markEscapable(r, c + 1, grid);
}

int numEnclaves(vector<vector<int>>& grid) {
    int rows = grid.size(), cols = grid[0].size();

    for (int r = 0; r < rows; r++) {
        markEscapable(r, 0, grid);
        markEscapable(r, cols - 1, grid);
    }
    for (int c = 0; c < cols; c++) {
        markEscapable(0, c, grid);
        markEscapable(rows - 1, c, grid);
    }

    int enclaves = 0;
    for (int r = 0; r < rows; r++) {
        for (int c = 0; c < cols; c++) {
            if (grid[r][c] == 1) enclaves++;
        }
    }
    return enclaves;
}

int main() {
    vector<vector<int>> grid = {
        {0, 0, 0, 0},
        {1, 0, 1, 0},
        {0, 1, 1, 0},
        {0, 0, 0, 0}
    };
    cout << numEnclaves(grid) << endl; // 3
    return 0;
}
\end{lstlisting}

\subsection*{Dry run}

\dryrun{Take the $4\times4$ grid above.}

Border land cells: $(1,0)$ (column $0$). $\textit{markEscapable}(1,0)$:
sets $(1,0)$ to $0$; its only land neighbor would need to be checked,
but $(0,0)=0,(2,0)=0,(1,1)=0$ -- no further spread. No other border
cell is land. The remaining land cells $(1,2),(2,1),(2,2)$ form a
connected interior region untouched by the border traversal -- they
are never adjacent to the border, and $(1,0)$'s region doesn't reach
them. Final count: $3$ enclaves.

\begin{complexitybox}
\textbf{Time Complexity.} $O(\textit{rows} \times \textit{cols})$.
\textbf{Space Complexity.} $O(\textit{rows} \times \textit{cols})$
recursion stack in the worst case.
\end{complexitybox}

\begin{takeawaybox}
Surrounded Regions and Number of Enclaves are, structurally, the same
algorithm: seed a traversal from the boundary, and treat whatever it
fails to reach as the answer. Once this boundary-first elimination
pattern is recognized, both problems reduce to writing the same dozen
lines of code with different variable names.
\end{takeawaybox}
